\label{sec:introducao}

O celular tornou-se uma porta de entrada para muitos serviços públicos brasileiros. Benefícios previdenciários, registros de vacinação, documentos e agendamentos de saúde podem ser consultados em aplicativos de órgãos federais, estaduais e municipais. Na coleta feita em setembro de 2026, por exemplo, o Meu INSS e o Meu SUS Digital apareciam na Google Play com mais de 50 milhões de instalações cada um.

Esses aplicativos podem conter kits de desenvolvimento de software (SDKs), componentes criados por outras empresas para oferecer funções prontas, como mapas, notificações e estatísticas de uso. Parte desses componentes coleta identificadores e informações sobre o uso do aplicativo, o que caracteriza o rastreamento por terceiros \cite{binns2018}. Esse rastreamento é comum em aplicativos comerciais \cite{binns2018,kollnig2021before} e também aparece no setor público: em 1.166 aplicativos governamentais de 71 países, Samarasinghe et~al.~\cite{samarasinghe2022} encontraram rastreadores do Google em 37\% deles e, em 27\%, o envio de informações sensíveis, como identificadores do usuário ou do dispositivo, a terceiros ou por conexões sem criptografia. Nos estudos examinados até agora, não há medição equivalente para o conjunto dos aplicativos governamentais brasileiros, nem a comparação entre os SDKs encontrados e o que cada aplicativo declara ao cidadão.

Esta pesquisa se concentra nesse ponto: identificar os SDKs de rastreamento presentes nos aplicativos oficiais e medir em que proporção a coleta e o compartilhamento de dados feitos por esses componentes estão declarados nos documentos publicados pelos próprios órgãos. A referência para essa medição não é a comparação com aplicativos comerciais, mas aquilo que cada serviço declara, examinado à luz da Lei Geral de Proteção de Dados Pessoais (LGPD). A dependência de fornecedores estrangeiros também será discutida, sem que a simples presença de um SDK seja tratada como transferência internacional de dados.

O objetivo geral é medir a proporção de pares aplicativo--SDK de rastreamento sem declaração compatível, isto é, de combinações entre um aplicativo e um SDK nele encontrado cuja coleta de dados não aparece nas declarações do aplicativo. Com base na literatura, a hipótese é que essa proporção chega a pelo menos 10\%. Para testá-la, o trabalho seguirá seis etapas: a revisão, um a um, dos 376 aplicativos já classificados como oficiais, pelo nome jurídico do desenvolvedor; o download da versão atual de cada um na Google Play, com o registro do hash do arquivo; a identificação dos SDKs por três detectores, cuja concordância será medida pelo $\kappa$ de Fleiss; a leitura da política de privacidade e da seção ``Segurança dos dados'' de cada aplicativo, para classificar cada par; a observação do tráfego de rede de uma subamostra executada em emulador; e a contagem dos pares sem declaração compatível, com o cálculo da sua proporção.

O estudo se limita ao Android, que concentra a maior parte dos celulares em uso no Brasil e cujos pacotes podem ser obtidos e analisados com ferramentas abertas, o que não ocorre no iOS. A seleção dos aplicativos e um piloto com 20 pacotes já foram realizados, assim como as buscas e a triagem inicial da revisão sistemática da literatura, cuja leitura completa está em andamento; a coleta completa e o teste da hipótese ocorrerão no Projeto Capstone II. Ao final, espera-se obter a proporção de pares sem declaração compatível, a lista das divergências encontradas em cada aplicativo, a medida de concordância entre os detectores e um conjunto de dados reprodutível.

\subsecao{Problema de pesquisa}\label{sec:problema}

Em que medida os SDKs de rastreamento de terceiros embarcados nos aplicativos Android oficiais do governo brasileiro estão declarados nas políticas de privacidade e na seção ``Segurança dos dados'' da Google Play?

Para responder a essa pergunta, a pesquisa separa duas evidências que não são equivalentes: o SDK identificado no código do pacote e o tráfego de fato observado durante a execução, que será examinado em uma subamostra. Encontrar um componente no código não basta para afirmar que ele enviou dados, e muito menos que houve infração à LGPD.

Os trabalhos brasileiros localizados estudam aplicativos específicos, como o Monitora Covid-19 \cite{lemos2022} e o Conecte SUS \cite{vaz2022}, ou analisam políticas de privacidade à luz da LGPD sem examinar o comportamento dos aplicativos \cite{reis2023}. A referência internacional mais próxima, Samarasinghe et~al.~\cite{samarasinghe2022}, não apresenta resultados separados para o Brasil nem compara detectores aplicados às mesmas versões dos pacotes. Essa lacuna é discutida na Seção~\ref{sec:ref-sintese}.

\subsecao{Justificativa}\label{sec:1.3}

Quem precisa consultar um benefício do Instituto Nacional do Seguro Social (INSS) ou acessar um serviço público de saúde nem sempre pode escolher outro fornecedor para realizar a mesma tarefa. Por isso, é especialmente importante saber quais empresas participam do funcionamento dos aplicativos oficiais e que informações elas podem receber. A relevância desse exame tem três dimensões: jurídica, de soberania de dados e metodológica.

Do ponto de vista jurídico, a LGPD define como sensível o dado referente à saúde e o dado biométrico (Art. 5º, II), restringe as hipóteses de tratamento desses dados (Art. 11), vincula o tratamento feito pelo poder público à sua finalidade pública (Art. 23) e limita a transferência internacional de dados pessoais a casos específicos (Art. 33) \cite{brasil2018}. Se um aplicativo de saúde publicado por um órgão público envia identificadores a um terceiro, a base legal desse envio precisa existir e ser demonstrável. Quando o envio não aparece na política publicada pelo próprio órgão, a divergência é um sinal de alerta, ainda que não prove, sozinha, a falta de base legal.

Quanto à soberania de dados, o rastreamento em aplicativos se concentra em poucas grandes empresas estrangeiras, como Google e Facebook \cite{binns2018,kollnig2021before}. Dados gerados pelo cidadão no uso de um serviço público podem, por uma dependência técnica assumida por padrão no desenvolvimento Android, ser processados sob outra jurisdição. A presença de SDKs de empresas estrangeiras já permite falar em dependência de fornecedores estrangeiros; afirmar que dados foram enviados ao exterior, porém, exige observar o tráfego, o que está previsto para uma subamostra (Seção~\ref{sec:trafego}). Medir a extensão dessa dependência é condição para discuti-la.

Há, por fim, uma razão metodológica. Duas ferramentas podem examinar o mesmo pacote e indicar quantidades diferentes de rastreadores. Antes de interpretar uma taxa como característica dos aplicativos, é preciso saber quanto ela depende do detector utilizado, e por isso o estudo compara detectores aplicados aos mesmos pacotes.

\subsecao{Objetivos}\label{sec:objetivos}

\subsubsecao{Objetivo geral}\label{sec:objetivo-geral}

Medir a proporção de pares aplicativo--SDK de rastreamento de terceiros sem declaração compatível nos aplicativos Android governamentais brasileiros, comparando os SDKs encontrados no código dos pacotes com a política de privacidade e a seção ``Segurança dos dados'' da Google Play, à luz da LGPD.

\subsubsecao{Objetivos específicos}\label{sec:1.4}

Para alcançar o objetivo geral, são propostos os seguintes objetivos específicos:

\begin{enumerate}[a)]
\item revisar, um a um, os 376 aplicativos classificados como oficiais, confirmando o órgão responsável e a esfera de governo pelo nome jurídico do desenvolvedor exibido na Google Play, e calcular a taxa de erro da classificação automática;
\item baixar a versão atual de cada aplicativo na Google Play, com o registro do hash SHA-256, da versão e da data, e listar os SDKs de rastreamento de cada pacote com três detectores baseados nas listas abertas Exodus Privacy e Disconnect;
\item ler a política de privacidade e a seção ``Segurança dos dados'' de cada aplicativo, classificar cada par aplicativo--SDK quanto ao que está declarado e calcular a proporção de pares sem declaração compatível;
\item executar uma subamostra de aplicativos em um emulador Android, registrar o tráfego de rede e identificar os domínios contatados, o país de destino e se esses destinos aparecem nas declarações;
\item calcular a concordância entre os três detectores pelo $\kappa$ de Fleiss e medir quanto a proporção de aplicativos com rastreador muda conforme o detector escolhido.
\end{enumerate}

\subsecao{Hipótese}\label{sec:1.5}

A pesquisa tem uma única hipótese, que será testada no Projeto Capstone II e se apoia em quatro resultados da literatura, discutidos na Seção~\ref{sec:revisao-bibliografica}. Andow et~al.~\cite{andow2020} confrontaram os fluxos de dados de 13.796 aplicativos Android com as respectivas políticas de privacidade, considerando também quem recebia cada dado, e verificaram que até 42,4\% dos aplicativos omitiam ou declaravam de forma incorreta seus fluxos de dados sensíveis. No setor público, Samarasinghe et~al.~\cite{samarasinghe2022} observaram, em 27\% de 1.166 aplicativos governamentais, o envio de informações sensíveis a terceiros ou por conexões sem criptografia, e Carlsson et~al.~\cite{carlsson2022}, ao comparar o tráfego de 32 aplicativos do setor público finlandês com suas políticas, encontraram divergências entre o que elas diziam e os dados enviados a serviços de análise de terceiros. Por fim, Kollnig et~al.~\cite{kollnig2022ios} apontaram, em 24 mil aplicativos Android e iOS, a persistente falta de transparência sobre o rastreamento. Esses estudos indicam que a ausência de declaração compatível aparece em parte relevante dos aplicativos analisados, inclusive nos públicos, o que leva à hipótese apresentada na Tabela~\ref{tab:hipotese}.

\begin{table}[htbp]
\centering
\small
\caption{Hipótese e critério de decisão}\label{tab:hipotese}
\begin{tabularx}{\textwidth}{|l|>{\raggedright\arraybackslash}X|>{\raggedright\arraybackslash}X|}
\hline
\textbf{Hipótese} & \textbf{Enunciado} & \textbf{Critério de decisão} \\
\hline
H1 & Ao menos 10\% dos pares aplicativo--SDK de rastreamento detectados nos aplicativos Android governamentais brasileiros não terão declaração compatível na política de privacidade nem na seção ``Segurança dos dados'' da Google Play, consideradas as categorias de dados e as finalidades. & A H1 é confirmada se a proporção de pares sem declaração compatível, entre os pares avaliáveis, for de pelo menos 10\%, e refutada se for menor. Pares sem documentos disponíveis ou com informação insuficiente são reportados à parte. \\ \hline
\end{tabularx}
\fonte{Elaborado pelo autor (2026), com base em Andow et~al.~\cite{andow2020}, Samarasinghe et~al.~\cite{samarasinghe2022}, Carlsson et~al.~\cite{carlsson2022} e Kollnig et~al.~\cite{kollnig2022ios}.}
\end{table}

O patamar de 10\% foi fixado bem abaixo das taxas da literatura porque os estudos citados usam outras populações e contam aplicativos, e não pares aplicativo--SDK; por isso, a hipótese é conservadora. Como serão examinados todos os aplicativos oficiais identificados, e não uma amostra, a decisão usa a proporção observada, e o intervalo de confiança de 95\% (IC95\%) indicará a precisão da estimativa, com ressalva na conclusão se incluir o patamar. Uma proporção abaixo de 10\% indicará ausência de declaração compatível menos frequente que o patamar desta pesquisa, sem permitir concluir que os aplicativos governamentais são mais transparentes do que as populações dos estudos anteriores. O patamar é um critério da pesquisa, e não uma tolerância prevista na LGPD.

As regras de classificação serão fixadas com o orientador antes da coleta principal, e o piloto não será usado como teste da hipótese. A concordância entre os detectores, objeto do quinto objetivo específico, não constitui uma segunda hipótese: trata-se de uma verificação do instrumento, que indica o quanto a lista de SDKs usada no teste depende da ferramenta escolhida (Seção~\ref{sec:deteccao}).

\subsecao{Estrutura do trabalho}\label{sec:estrutura}

Este trabalho está organizado em cinco seções. Após esta introdução, a Seção~\ref{sec:revisao-bibliografica} reúne a fundamentação teórica e os trabalhos correlatos, incluindo os estudos com proposta de inclusão na revisão sistemática. A Seção~\ref{sec:metodologia} descreve a metodologia: as etapas do estudo, o protocolo da revisão sistemática, os aspectos éticos, os recursos e o cronograma. A Seção~\ref{sec:resultados} apresenta os resultados preliminares obtidos na preparação do estudo, os resultados esperados e as limitações, e a Seção~\ref{sec:conclusao} traz as considerações finais e os próximos passos no Projeto Capstone II.
